\documentclass[../DoAn.tex]{subfiles}
\begin{document}
\section{Kết luận}
Đồ án đã xây dựng một hệ thống hỏi đáp và rà soát văn bản pháp luật Việt Nam trên dữ liệu Bộ Pháp điển, kết hợp RAG, đồ thị tri thức và tác tử AI. So với các cổng tra cứu chính thống, hệ thống hiểu được câu hỏi bằng ngôn ngữ đời thường và tổng hợp câu trả lời thay vì chỉ trả về danh sách văn bản. So với các thư viện pháp luật thương mại, hệ thống trả lời được những câu hỏi chưa từng được đăng và rà soát được hợp đồng do người dùng tải lên. So với các trợ lý AI đa năng, mỗi câu trả lời và cảnh báo của hệ thống đều gắn với điều luật thực sự tra được, người dùng có thể mở toàn văn để tự kiểm tra, và hệ thống chủ động kiểm tra quy định hiện hành với các câu hỏi về mức phạt, mức tiền.

Trong quá trình thực hiện, đồ án đã hoàn thành các chức năng chính đề ra ở Chương \ref{chapter:Introduction}: hỏi đáp và phân tích tình huống có trích dẫn; rà soát văn bản từ file, đường dẫn hoặc nội dung dán vào, có lưu lịch sử; bộ nhớ dài hạn, quản lý hội thoại, hỏi và nghe bằng giọng nói, đăng nhập và đồng bộ dữ liệu ẩn danh; công cụ quản trị cập nhật văn bản và theo dõi hiệu lực; văn phòng luật 3D. Hệ thống gồm khoảng 22 nghìn dòng mã nguồn, 299 kiểm thử đơn vị và ba bộ đánh giá tự động. Các đóng góp nổi bật gồm tìm kiếm lai cho văn bản pháp lý tiếng Việt nâng nDCG@10 từ 0,417 lên 0,632, luật buộc kiểm tra web giúp tác tử gọi đúng công cụ bắt buộc ở 97,6\% trường hợp, cơ chế kiểm chứng căn cứ cho rà soát văn bản và bộ nhớ dài hạn bảo vệ quyền riêng tư.

Đồ án còn một số hạn chế. Thứ nhất, các bộ đánh giá có quy mô nhỏ (50 câu hỏi truy xuất, 23 tình huống gọi công cụ, 14 kịch bản bộ nhớ) và dùng LLM làm người chấm, nên kết quả cần được kiểm chứng thêm bởi chuyên gia pháp lý trên tập dữ liệu lớn hơn. Thứ hai, tính năng rà soát văn bản mới được thử nghiệm trên một số hợp đồng mẫu, kết quả còn dao động giữa các lần chạy, và hệ thống chỉ bảo đảm căn cứ là điều luật có thật chứ chưa bảo đảm lập luận áp dụng luôn đúng. Thứ ba, mô hình dữ liệu hiện lưu một trạng thái hiệu lực cho mỗi điều, chưa lưu được nhiều phiên bản nội dung theo thời gian. Thứ tư, đồ thị tri thức hiện mới khai thác quan hệ liên quan có sẵn trong Pháp Điển, quan hệ tham chiếu trích từ nội dung còn thưa. Thứ năm, hệ thống phụ thuộc vào API của OpenAI và Tavily, phát sinh chi phí theo lượt dùng.

Bài học lớn nhất rút ra là với ứng dụng pháp lý, độ tin cậy phải được bảo đảm bằng các quy tắc kiểm tra trong mã nguồn chứ không thể chỉ dựa vào lời nhắc cho mô hình. Mọi lỗi nghiêm trọng gặp trong quá trình phát triển, từ đọc nhầm mã sắp xếp thành số hiệu văn bản, không chịu tra web khi đã có kết quả ``có vẻ khớp'', tới chép sai mã điều luật dài, đều được khắc phục bằng cách đưa quyết định quan trọng ra khỏi mô hình và kiểm chứng đầu ra của mô hình. Bài học thứ hai là cần đo trước khi tối ưu: bộ đánh giá không cần nhãn thủ công cho phép so sánh các cấu hình truy xuất một cách khách quan, nhờ đó phát hiện cấu hình ba nguồn tuy phức tạp hơn nhưng không tốt hơn cấu hình hai nguồn.

\section{Hướng phát triển}
Trước hết, cần hoàn thiện các chức năng hiện có. Bộ đánh giá nên được mở rộng lên vài trăm câu hỏi có nhãn do người có chuyên môn pháp lý kiểm tra, đồng thời xây dựng bộ hợp đồng mẫu có danh sách điều khoản trái luật đã biết để đo độ phủ và độ chính xác của tính năng rà soát. Tính năng rà soát nên tra thêm một truy vấn theo chính nội dung điều khoản để giảm dao động, gộp các vấn đề cùng một điều khoản thành một cảnh báo và đọc ngày ký hợp đồng để đối chiếu với luật có hiệu lực tại thời điểm đó. Bộ trích xuất cần tích hợp nhận dạng ký tự quang học cho văn bản quét.

Hướng phát triển quan trọng nhất là lưu trữ điều luật theo thời gian. Thiết kế đề xuất tách danh tính của điều luật khỏi nội dung của nó: mỗi điều có một mã cố định và nhiều phiên bản nội dung, mỗi phiên bản gắn với khoảng hiệu lực và văn bản làm phát sinh phiên bản đó; quan hệ thay thế, sửa đổi, bãi bỏ, bổ sung được lưu ở mức điều kèm phạm vi (cả điều, một khoản hay một điểm) và ngày có hiệu lực. Mọi truy vấn khi đó nhận thêm một ngày cần xét, mặc định là hôm nay hoặc ngày ký hợp đồng. Với văn bản hợp nhất, vốn theo Pháp lệnh số 01/2012/UBTVQH13 \cite{plhn2012} không đặt ra quy phạm mới, mỗi điều trong văn bản hợp nhất được ghép với điều gốc và trở thành một phiên bản nội dung mới của điều gốc, với ngày hiệu lực lấy từ văn bản sửa đổi được nêu trong chú thích; vì vậy chú thích cuối trang cần được tách ra lưu riêng thay vì loại bỏ như hiện nay. Thiết kế này bảo đảm điều luật cũ không bao giờ bị mất, điều luật mới được ưu tiên đúng thời điểm và việc tra lại một hợp đồng cũ vẫn ra đúng luật áp dụng lúc ký, phù hợp với nguyên tắc hiệu lực của Luật Ban hành văn bản quy phạm pháp luật \cite{luatbhvb2015}.

Về dài hạn, hệ thống có thể được mở rộng theo ba hướng. Hướng thứ nhất là làm giàu đồ thị tri thức bằng bộ phân tích trích dẫn ``Điều X văn bản Y'' trong nội dung điều luật để tạo đầy đủ quan hệ tham chiếu, và áp dụng tóm tắt theo cộng đồng của Graph RAG \cite{edge2024graphrag} cho các câu hỏi cần cái nhìn tổng quát. Hướng thứ hai là giảm phụ thuộc vào dịch vụ ngoài bằng mô hình nhúng và mô hình ngôn ngữ mã nguồn mở được tinh chỉnh trên văn bản pháp luật tiếng Việt, đặc biệt cho các tác vụ phụ như viết lại câu hỏi và trích xuất bộ nhớ. Hướng thứ ba là mở hệ thống ra các kênh khác: đóng gói các công cụ tra cứu Pháp Điển và rà soát văn bản thành một máy chủ theo giao thức Model Context Protocol để các trợ lý AI khác sử dụng, tích hợp bản đồ để chỉ dẫn toà án, văn phòng công chứng gần nhất, và tích hợp lịch để nhắc các thời hạn trong hợp đồng vừa rà soát.

\end{document}
